% !TeX root = ../main.tex
% Rebuilt from: continuum-mechanics.tex, ch02-stress-and-intact-rock.tex
\chapter{Stress, Strain and Intact Rock Behaviour}
\label{ch:rock-mechanics}

The Kirsch solution~\citep{kirsch1898}, the convergence--confinement method and every finite
element analysis in this book rest on the same small set of continuum
relations. Setting them out once, at the start, avoids restating their
assumptions at every later application --- and makes visible the fact that
those assumptions are shared. A closed-form stress distribution and a
sophisticated numerical model can be wrong for the same reason.

The chapter then turns from the relations that govern stress to the
criteria that limit it. It sets out the Mohr--Coulomb, Hoek--Brown and
Griffith descriptions of intact rock strength, distinguishes the stress the
ground imposes from the stress the rock can sustain, and derives the
strength parameters of four intact rocks that
Chapter~\ref{ch:rock-mass-parameters} carries forward to the rock mass.

\section{Continuum description of stress and strain}
\label{sec:continuum-description-of-stress-and-strain}

\subsection{Stress}
\label{sec:stress}

\paragraph{Traction and the stress tensor.}
Consider a plane through a point in a loaded body. The material on one side
exerts a force on the material on the other, and the traction is that force
per unit area. Its magnitude and direction depend on the orientation of the
plane, so a single vector cannot describe the state of stress at a point;
nine components are required, three tractions on each of three mutually
perpendicular planes, forming the stress tensor
\[
\sigma_{ij} =
\begin{bmatrix}
\sigma_{xx} & \tau_{xy} & \tau_{xz}\\
\tau_{yx} & \sigma_{yy} & \tau_{yz}\\
\tau_{zx} & \tau_{zy} & \sigma_{zz}
\end{bmatrix}.
\]
Moment equilibrium of an infinitesimal element requires
$\tau_{ij} = \tau_{ji}$, so only six components are independent.

Compression is taken as positive throughout, following rock mechanics
practice. This is the opposite of the structural convention and is the single
most common source of sign errors where the two disciplines meet.

\paragraph{Principal stresses.}
For any state of stress there exists an orientation on which the shear
components vanish and only normal stresses act. These are the principal
stresses, obtained as the eigenvalues of $\sigma_{ij}$, with the principal
directions as its eigenvectors. They are ordered
$\sone \geq \stwo \geq \sthree$.

Principal stresses matter because rock strength is expressed in terms of them.
Every failure criterion in this chapter is written as a relationship between
$\sone$ and $\sthree$, and applying one therefore requires the stress state to
be resolved into principal components first.

\paragraph{Invariants and the deviatoric split.}
The principal values are properties of the stress state, not of the axes used
to describe it, so combinations of the components that are independent of
orientation --- the invariants --- carry physical meaning. The first is the
sum of the normal stresses, from which the mean stress
$p = (\sone+\stwo+\sthree)/3$ follows.

Separating the mean from the remainder gives the hydrostatic and deviatoric
parts. The hydrostatic part changes volume; the deviatoric part changes shape
and drives shear failure. The split is fundamental to rock behaviour because
rock strength depends on both: the deviatoric stress causes failure and the
mean stress determines the resistance available. That dependence on confining
pressure distinguishes rock from metals, whose yield is essentially
pressure-independent, and it is why rock criteria rise with confinement in
$\sone$--$\sthree$ space, whether as a straight line (Mohr--Coulomb) or as a
curve (Hoek--Brown), rather than as the unit-slope line that
pressure-independent yield would give.

\subsection{Equilibrium}
\label{sec:equilibrium}

\paragraph{The equilibrium conditions.}
Any admissible stress field must satisfy equilibrium at every point. In
Cartesian coordinates with body force $f_i$,
\[
\frac{\partial \sigma_{ij}}{\partial x_j} + f_i = 0 ,
\]
and for a circular opening the cylindrical form is more useful. Neglecting
body force and assuming axial symmetry, radial equilibrium reduces to
\[
\frac{\mathrm{d}\sigma_r}{\mathrm{d}r} + \frac{\sigma_r-\sigma_\theta}{r} = 0 ,
\]
which is the starting point for both the Kirsch solution
(Section~\ref{sec:kirsch-appendix}) and the elasto-plastic
ground reaction curve of Chapter~\ref{ch:ccm}. Any candidate stress
distribution around an opening must satisfy it; a numerical result that does
not has failed a check that costs nothing to make.

\subsection{Dimensional idealisation}
\label{sec:plane-idealisations}

\paragraph{Plane strain, plane stress and axisymmetry.}
Three-dimensional analysis is rarely necessary and often unaffordable, and
three idealisations make most underground problems tractable.

\emph{Plane strain} assumes no strain along one axis. A long tunnel remote
from the face satisfies this well, since the ground cannot displace
longitudinally, and it is the standard idealisation for tunnel cross-sections.
\emph{Plane stress} assumes no stress on one pair of faces and applies to thin
plates; it is rare underground. \emph{Axisymmetry} exploits rotational
symmetry about an axis and suits a circular tunnel in a hydrostatic field, or
a shaft.

The important limitation concerns the face. Near it the ground does displace
longitudinally, so the plane strain assumption fails exactly where support is
installed and where most of the convergence occurs. Two-dimensional analysis
handles this through a stress relaxation factor that releases only part of the
in-situ stress before support is activated --- a device, not a solution, and
the reason three-dimensional analysis is required for junctions, portals and
sequenced excavations (Chapter~\ref{ch:numerical-modelling}).

\section{Constitutive behaviour}
\label{sec:constitutive-behaviour}

\subsection{Strain and compatibility}
\label{sec:strain-and-compatibility}

\paragraph{The strain tensor.}
Strain measures relative displacement. For small displacements the
infinitesimal strain tensor is
\[
\varepsilon_{ij} = \tfrac{1}{2}
\left( \frac{\partial u_i}{\partial x_j} + \frac{\partial u_j}{\partial x_i} \right),
\]
symmetric like the stress tensor and, like it, decomposable into volumetric
and deviatoric parts.

One notational trap recurs. Engineering shear strain $\gamma_{ij}$ is twice
the tensor shear strain $\varepsilon_{ij}$. Both appear in the literature,
frequently without statement, and confusing them introduces a factor of two
into the shear modulus.

\paragraph{Compatibility.}
Six strain components are derived from three displacement components, so they
cannot be independent. The compatibility conditions express that constraint: a
strain field satisfying them corresponds to a continuous displacement field,
and one that does not implies gaps or overlaps in the material.

For closed-form solutions this is a working requirement rather than a formal
nicety, because a solution is constructed by satisfying equilibrium,
compatibility and the boundary conditions simultaneously. In numerical
analysis compatibility is enforced automatically by the displacement
formulation, which is one reason finite element methods handle complex
geometry so readily.

\subsection{Constitutive relations}
\label{sec:constitutive-relations}

\paragraph{Isotropic elasticity.}
Equilibrium and compatibility are insufficient to determine stress and strain:
a relation between them is required. For a linear elastic isotropic material,
generalised Hooke's law gives
\[
\varepsilon_{xx} = \tfrac{1}{E}\left[\sigma_{xx} - \nu(\sigma_{yy}+\sigma_{zz})\right],
\qquad
\gamma_{xy} = \frac{\tau_{xy}}{G},
\]
with corresponding expressions for the remaining components. Only two
constants are independent, and the others follow:
$G = E/[2(1+\nu)]$ and $K = E/[3(1-2\nu)]$.

Under plane strain the relations reduce with $\varepsilon_{zz}=0$, giving
$\sigma_{zz} = \nu(\sigma_{xx}+\sigma_{yy})$: the out-of-plane stress is not
zero but is determined by the in-plane state. Overlooking this term is a
common error when interpreting two-dimensional output, since it is often the
intermediate principal stress.

\paragraph{Anisotropy.}
Bedded and foliated rocks are not isotropic. Their stiffness parallel to
bedding differs from that perpendicular to it, sometimes by a factor of two or
more, and the same holds for strength. Transverse isotropy --- one axis of
rotational symmetry normal to the bedding --- describes this adequately and
requires five independent elastic constants rather than two.

The obstacle is measurement. Determining five constants requires specimens
cored at several orientations to the fabric, which multiplies the testing
programme and demands core that survives being taken at unfavourable angles.
Anisotropy is therefore frequently acknowledged in the geological description
and ignored in the analysis. Where the fabric is strong and unfavourably
oriented relative to the excavation, that omission is not conservative: it
misplaces both the stiffness and the failure plane.

\subsection{Formulations for circular openings}
\label{sec:formulations-for-circular-openings}

\paragraph{Cylindrical coordinates.}
Circular openings are naturally described in $r$--$\theta$ coordinates, in
which stress has radial, tangential and shear components and the equilibrium
equation takes the reduced form given above. Two features of the solution
recur throughout the book.

The tangential stress $\sigma_\theta$ at the boundary is the quantity that
governs. It concentrates there, and where it exceeds the rock strength,
spalling follows. The radial stress $\sigma_r$ falls to the support pressure
at the boundary and to zero if the opening is unsupported, which is why
excavation reduces confinement and drives the rock towards a weaker part of
its failure envelope even though the far-field stress is unchanged.

\paragraph{Strain energy.}
An elastically stressed rock mass stores energy, with density
$U = \tfrac{1}{2}\sigma_{ij}\varepsilon_{ij}$. Excavation releases some of it
--- part absorbed in deformation and in work done on the support, the
remainder available to drive fracture.

At moderate depth the released energy dissipates quietly. At high stress in
brittle rock it need not: where the release exceeds what the rock and support
can absorb, failure becomes violent. This is the basis of the energy treatment
of rockburst in Chapter~\ref{ch:high-stress}, and it is the reason depth changes
not only the magnitude of the design problem but its character.

\section{Failure theory and criteria}
\label{sec:failure-theory-and-criteria}

\subsection{Inelastic behaviour and graphical methods}
\label{sec:inelastic-behaviour-and-graphical-methods}

\paragraph{Yield, flow and dilation.}
Beyond the elastic limit, further vocabulary is needed. The \emph{yield
surface} bounds the elastic region in stress space; the failure criteria in
this chapter are yield surfaces. The \emph{plastic potential} governs the
direction of plastic straining. Where the two coincide the flow rule is
\emph{associated}; where they do not it is \emph{non-associated}.

For rock the distinction is practical rather than academic. An associated flow
rule with a Mohr--Coulomb yield surface predicts dilation at an angle equal to
the friction angle, typically 30 to 40 degrees, which grossly overstates the
volume increase actually observed. Non-associated flow with a dilation angle
well below the friction angle --- often a third of it, and approaching zero at
high confinement --- is the realistic choice, and dilation angle is a required
input to most numerical models.

Post-peak behaviour completes the picture. \emph{Perfectly plastic} response
maintains strength after yield; \emph{strain-softening} loses it
progressively; \emph{brittle} response drops it abruptly. Rock mass quality
largely determines which applies, with strong massive rock behaving brittly
and weak jointed rock nearly plastically. The choice changes the extent of the
computed yield zone substantially, and it is an assumption rather than a
measurement.

\paragraph{The Mohr circle.}
Plotting normal against shear stress represents a two-dimensional stress state
as a circle through $(\sone,0)$ and $(\sthree,0)$, on which each point gives
the stress on a plane of corresponding orientation. Failure occurs where the
circle touches the strength envelope, and the tangent point identifies the
plane on which it occurs.

For underground work the construction earns its place through discontinuities.
Given the stress state and a joint of known orientation, the Mohr circle
returns the normal and shear stress acting on that joint directly, and
comparison against its shear strength (Chapter~\ref{ch:engineering-geology})
determines whether it will slip. Three-dimensional states plot as three
circles, and the region between them contains the admissible stress
combinations --- a reminder that $\stwo$, ignored by both Mohr--Coulomb and
Hoek--Brown, is nonetheless present and does affect strength.

\subsection{Families of failure criteria}
\label{sec:material-failure-theory}

A failure criterion states the combination of stresses at which a material
ceases to carry further load in its existing form. No single criterion holds
for all materials, and none is exact for any. Each is a simplified model
developed for a class of material and a range of stress, so the choice among
them is a design decision that has to be justified against the problem in
hand. Three families are distinguished by what they take failure to be.

\paragraph{Stress-based criteria.}
Phenomenological criteria define failure as the attainment of a critical
stress condition, without reference to the mechanism that produces it.
Table~\ref{tab:phenomenological_criteria} lists those in common use. Two of
them carry almost all of rock engineering practice. The Hoek--Brown criterion
represents the curved, confinement-dependent strength of rock, and the
Mohr--Coulomb criterion supplies the cohesion and friction angle in which
strength enters most design calculations and many numerical models. The usual
sequence is therefore to derive rock strength with Hoek--Brown and to convert
it to equivalent Mohr--Coulomb parameters where a calculation requires them.

\begin{table}[htbp]
\caption{Common phenomenological failure criteria.}
\label{tab:phenomenological_criteria}
\centering
\small
\begin{tabularx}{\textwidth}{l Y Y}
\toprule
\textbf{Criterion} & \textbf{Main idea} & \textbf{Typical use} \\
\midrule
Mohr--Coulomb & Linear cohesion--friction envelope & Soils, weak rocks, rock masses, design calculations, numerical modelling \\
Rankine & Maximum principal stress criterion & Brittle tensile failure and simple stress-limit checks \\
Tresca & Maximum shear stress criterion & Ductile metals under simple loading \\
von Mises & Distortional strain energy criterion & Ductile metals and plasticity models \\
Hoek--Brown & Nonlinear confinement-dependent rock strength & Intact rock and rock mass strength estimation \\
Drucker--Prager & Smooth pressure-dependent approximation of Mohr--Coulomb & Numerical modelling and plasticity formulations \\
Cam--Clay & Critical-state soil mechanics model & Normally consolidated clays and advanced soil modelling \\
\bottomrule
\end{tabularx}
\end{table}

\paragraph{Fracture mechanics and energy methods.}
Fracture-based criteria address the initiation and propagation of cracks
rather than the overall stress state. Griffith theory, the earliest of them,
holds that a crack extends when the energy released by its growth is
sufficient to create the new surface~\citep{griffith1921rupture}. The approach
suits brittle materials such as rock, concrete and glass, and rock in
particular, because rock is a flawed material: failure initiates at
microcracks, bedding planes, joints and other defects, and not from a
homogeneous stress state.

Energy-based methods extend the same reasoning to conditions in which a
stress criterion alone is insufficient. These include anisotropic materials,
complex loading paths, dynamic loading, fracture propagation and advanced
numerical modelling. The energy treatment of rockburst in
Chapter~\ref{ch:high-stress} is the principal instance in this book.

\paragraph{Selecting a criterion.}
The criterion should follow from the engineering problem and not from
convenience. Table~\ref{tab:criterion_selection} sets out the considerations:
the material, the scale, the level of confinement, the expected mechanism, the
data available, the form the analysis accepts and the purpose of the
calculation. To adopt Mohr--Coulomb because it is simple justifies nothing. To
adopt it because the analysis requires a linear strength envelope over the
confinement range of the project, with parameters derived from Hoek--Brown
over that range, is a statement that can be examined and defended.

\begin{table}[htbp]
\caption{Considerations in selecting a failure criterion.}
\label{tab:criterion_selection}
\centering
\small
\begin{tabularx}{\textwidth}{l Y Y Y}
\toprule
\textbf{Consideration} & \textbf{Governing question} & \textbf{Suitable criterion} & \textbf{Treatment in this book} \\
\midrule
Material type
& Is the material soil, intact rock, jointed rock, or concrete?
& Rock: HB/MC; soil: MC/DP/Cam--Clay; concrete: MC or fracture-based checks
& Rock projects use HB first, then MC when required. \\

Scale
& Is the analysis at the scale of intact laboratory samples or of the rock mass around a tunnel?
& Intact scale: intact HB; rock mass scale: generalised HB or equivalent MC
& This chapter uses intact HB; Chapter~\ref{ch:rock-mass-parameters} uses generalised HB for rock masses. \\

Confining stress
& Is confinement low, moderate, or high?
& Low confinement: MC may be adequate; high confinement: HB preferred
& Deep tunnels and caverns need HB sensitivity checks. \\

Failure mechanism
& Is failure brittle, shear-dominated, structurally controlled, or squeezing?
& Brittle: HB/Griffith concepts; shear: MC; structurally controlled: discontinuity and empirical methods
& Chapter~\ref{ch:ground-behaviour} uses failure mechanism identification before support design. \\

Available data
& Are laboratory strengths, GSI and RMR or Q available, or only a limited description?
& Detailed data: HB or advanced models; limited data: MC or empirical correlations
& Practice moves from laboratory data to classification and then to design parameters. \\

Analysis requirement
& What criterion does the intended analysis accept?
& Continuum codes generally accept both; several closed-form solutions require Mohr--Coulomb
& Use MC where specified, but justify how it was derived. \\

Purpose
& Is the purpose preliminary screening, parameter derivation, or final design?
& Preliminary: MC/empirical; parameter derivation: HB; detailed design: calibrated numerical model
& This chapter derives parameters; Chapters~\ref{ch:empirical-design} to~\ref{ch:numerical-modelling} compare empirical, analytical and numerical methods. \\
\bottomrule
\end{tabularx}
\end{table}

For rock the sequence adopted in this book has four steps. The intact
Hoek--Brown criterion gives the confined compressive strength and the tensile
strength from \(\sigma_{ci}\) and \(m_i\). The Hoek--Brown envelope is
converted to equivalent Mohr--Coulomb parameters \(c\) and \(\phi\) over a
stress range \(0 \leq \sigma_3 \leq \sigma_{3,\mathrm{fit,max}}\) specific to
the project. The condition of the rock mass is described by RMR, Q and GSI
(Chapter~\ref{ch:classification}). Finally, the generalised Hoek--Brown
criterion~\citep{hoek2019hb}, with GSI and the disturbance factor \(D\), gives
the strength of the rock mass, which is converted in turn to equivalent
Mohr--Coulomb parameters (Chapter~\ref{ch:rock-mass-parameters}). The sequence
is defensible because Hoek--Brown captures the nonlinear behaviour of rock,
while Mohr--Coulomb provides the \(c\) and \(\phi\) that many design
calculations require.

Figure~\ref{fig:failure_workflow} places the selection within the wider
design sequence. The present chapter occupies the middle of it, the selection
of a criterion for intact rock and the derivation of its parameters from
laboratory data; Chapters~\ref{ch:classification}
and~\ref{ch:rock-mass-parameters} carry the same logic from intact rock to the
jointed rock mass.

\begin{figure}[htbp]
\centering
\resizebox{\textwidth}{!}{%
\begin{tikzpicture}[node distance=0.45cm, font=\small]

\node[workflowbox, wfinput] (step1)
{\textbf{1. Understand}\\ material and site\\
\footnotesize geology, stress};

\node[workflowbox, wfstep, right=of step1] (step2)
{\textbf{2. Identify}\\ likely behaviour\\
\footnotesize brittle, ductile, jointed, squeezing};

\node[workflowbox, wfstep, right=of step2] (step3)
{\textbf{3. Select}\\ failure criterion\\
\footnotesize HB, MC, DP, Griffith, etc.};

\node[workflowbox, wfstep, right=of step3] (step4)
{\textbf{4. Derive}\\ parameters\\
\footnotesize \(c,\phi,\sigma_t,\sigma_{ci},m_i\)};

\node[workflowbox, wfcheck, right=of step4] (step5)
{\textbf{5. Analyse}\\ and validate\\
\footnotesize against laboratory and field data};

\draw[workflowarrow] (step1) -- (step2);
\draw[workflowarrow] (step2) -- (step3);
\draw[workflowarrow] (step3) -- (step4);
\draw[workflowarrow] (step4) -- (step5);

\draw[feedbackarrow] (step5.south) -- ++(0,-0.65) -| node[pos=0.25,below,font=\footnotesize]
{iterate when assumptions or results are inconsistent} (step1.south);

\end{tikzpicture}
}
\caption{The place of the failure criterion in the design sequence, from the description of the ground to the validation of the analysis.}
\label{fig:failure_workflow}
\end{figure}

\subsection{In-situ stress as the demand}
\label{sec:stress-and-the-in-situ-stress-field}

A failure criterion supplies a capacity. The demand against which it is
compared is the in-situ stress, whose origin, variation with depth and
measurement are the subject of Chapter~\ref{ch:in-situ-stress}. Only what is
needed to apply a criterion is restated here.

\paragraph{Vertical and horizontal components.}
In-situ stress is the natural stress in the ground before excavation. It is
generated by the weight of the overlying rock or soil, by tectonic plate
movements, by the geological history of erosion, uplift and residual stress,
and by local structures such as faults, folds and bedding. It is a
three-dimensional state, and the horizontal components may exceed the
vertical.

The vertical stress is estimated from the overburden,
\[
\sigma_v = \gamma z ,
\]
where \(\gamma\) is the unit weight of the overlying material and \(z\) the
depth below ground surface. The examples of this chapter take
\(\gamma = 27~\text{kN/m}^3 = 0.027~\text{MPa/m}\), so that
\(\sigma_v = 0.027z\) in megapascals and, at \(z=2000~\text{m}\),
\(\sigma_v = 0.027 \times 2000 = 54~\text{MPa}\). Unit weights by rock type,
which span 22 to 30~kN/m\textsuperscript{3}, are in
Table~\ref{tab:intact-props} of Appendix~\ref{app:constants}.

The horizontal stresses are expressed through in-situ stress ratios,
\[
\sigma_H = k_H \sigma_v ,
\qquad
\sigma_h = k_h \sigma_v ,
\qquad
\sigma_H \geq \sigma_h ,
\]
in which the uppercase subscript denotes the maximum horizontal stress and the
lowercase the minimum. For ground deforming elastically under gravity with no
lateral strain, the ratio is \(k = \nu/(1-\nu)\), which for \(\nu = 0.25\)
gives \(k = 0.25/(1-0.25) = 0.33\). The expression takes no account of
tectonic stress, stress history, erosion, faulting or locked-in residual
stress, and Section~\ref{sec:stress-components} shows it to be a lower bound.
Measured or project-specific ratios are preferred wherever they are
available.

\paragraph{Stress regimes.}
The relative magnitudes of \(\sigma_v\), \(\sigma_H\) and \(\sigma_h\) define
three tectonic regimes (Table~\ref{tab:stress_regimes}). The major principal
stress may therefore be vertical or horizontal, and \(\sigma_1=\sigma_v\)
cannot be assumed. The three components are estimated first and then ranked to
identify \(\sigma_1\), \(\sigma_2\) and \(\sigma_3\).

\begin{table}[htbp]
\caption{Common stress regimes in rock mechanics.}
\label{tab:stress_regimes}
\centering
\small
\begin{tabularx}{\textwidth}{l c Y}
\toprule
\textbf{Stress regime} & \textbf{Principal stress order} & \textbf{Typical interpretation} \\
\midrule
Normal faulting &
\(\sigma_v = \sigma_1 > \sigma_H = \sigma_2 > \sigma_h = \sigma_3\) &
Vertical stress is the largest stress; common in extensional settings. \\
Strike-slip faulting &
\(\sigma_H = \sigma_1 > \sigma_v = \sigma_2 > \sigma_h = \sigma_3\) &
Maximum horizontal stress is largest; vertical stress is intermediate. \\
Reverse faulting &
\(\sigma_H = \sigma_1 > \sigma_h = \sigma_2 > \sigma_v = \sigma_3\) &
Both horizontal stresses exceed vertical stress; common in compressional settings. \\
\bottomrule
\end{tabularx}
\end{table}

\paragraph{From depth to principal stresses.}
Figure~\ref{fig:insitu_stress_workflow} sets out the sequence from the depth
of the project to the comparison of demand with capacity. Where a project
supplies a single stress ratio \(k\), as the data sets of this chapter and of
Chapter~\ref{ch:rock-mass-parameters} do, the horizontal stress is
\(\sigma_{h,k} = k\sigma_v\) and the ranking reduces to a comparison of two
stresses,
\[
\sigma_1^{IS} = \max(\sigma_v,\sigma_{h,k}) ,
\qquad
\sigma_3^{IS} = \min(\sigma_v,\sigma_{h,k}) ,
\]
in which the superscript marks a principal stress acting in the ground before
excavation.

\begin{figure}[htbp]
\centering
\resizebox{\textwidth}{!}{%
\begin{tikzpicture}[node distance=0.27cm, font=\footnotesize]

\node[tinyworkflowbox, wfinput] (a)
{\textbf{1. Project depth}\\
determine \(z\) and \(\gamma\)};

\node[tinyworkflowbox, wfstep, right=of a] (b)
{\textbf{2. Vertical stress}\\
\(\sigma_v=\gamma z\)};

\node[tinyworkflowbox, wfstep, right=of b] (c)
{\textbf{3. Horizontal stresses}\\
\(\sigma_H=k_H\sigma_v\)\\
\(\sigma_h=k_h\sigma_v\)};

\node[tinyworkflowbox, wfstep, right=of c] (d)
{\textbf{4. Rank stresses}\\
identify \(\sigma_1^{IS}\), \(\sigma_2^{IS}\), \(\sigma_3^{IS}\)};

\node[tinyworkflowbox, wfcheck, right=of d] (e)
{\textbf{5. Use criterion}\\
compare stress demand\\
with strength capacity};

\draw[workflowarrow] (a) -- (b);
\draw[workflowarrow] (b) -- (c);
\draw[workflowarrow] (c) -- (d);
\draw[workflowarrow] (d) -- (e);

\end{tikzpicture}
}
\caption{Sequence from project depth and stress ratios to the principal stresses used in a failure assessment.}
\label{fig:insitu_stress_workflow}
\end{figure}

\section{Failure criteria for intact rock}
\label{sec:failure-criteria-for-intact-rock}

\subsection{The Mohr--Coulomb criterion}
\label{sec:mohr-coulomb-criterion}

\paragraph{The linear envelope.}
The Mohr--Coulomb criterion expresses the shear strength \(\tau\) on a plane
as a linear function of the normal stress \(\sigma_n\) acting on it,
\[
\tau = c + \sigma_n \tan\phi ,
\]
where \(c\) is the cohesion and \(\phi\) the friction angle. In terms of the
principal stresses \((\sigma_1 \geq \sigma_2 \geq \sigma_3)\), failure occurs
when
\[
\sigma_1 =
\frac{2c\cos\phi + \sigma_3(1+\sin\phi)}
{1-\sin\phi} ,
\]
which is a straight line in \(\sigma_1\)--\(\sigma_3\) space,
\[
\sigma_1 = \sigma_c + k_\phi\,\sigma_3 ,
\qquad
\sigma_c =
\frac{2c\cos\phi}{1-\sin\phi} ,
\qquad
k_\phi =
\frac{1+\sin\phi}{1-\sin\phi} .
\]
The cohesion mainly controls strength at low confinement, and the friction
angle controls how rapidly strength rises with confining stress.

\paragraph{Uniaxial and tensile strengths.}
The intercept \(\sigma_c\) is the uniaxial compressive strength predicted by
the criterion, obtained with \(\sigma_3 = 0\). Setting \(\sigma_1 = 0\)
instead gives the uniaxial tensile strength,
\[
\sigma_t =
-\frac{2c\cos\phi}{1+\sin\phi}
=
-\frac{\sigma_c(1-\sin\phi)}{1+\sin\phi} ,
\]
negative under the compression-positive convention.

\paragraph{Orientation of the failure plane.}
The failure plane is inclined to the direction of \(\sigma_3\), that is to the
plane on which \(\sigma_1\) acts, at
\[
\theta = 45^\circ + \frac{\phi}{2},
\]
and therefore at \(45^\circ - \phi/2\) to the direction of \(\sigma_1\)
itself. For \(\phi = 30^\circ\) the plane lies \(30^\circ\) from the
\(\sigma_1\) direction.

\paragraph{Mohr circle representation.}
A triaxial stress state plots as a Mohr circle of centre \(C\) and radius
\(R\),
\[
C = \frac{\sigma_1+\sigma_3}{2},
\qquad
R = \frac{\sigma_1-\sigma_3}{2} ,
\]
and failure occurs when the circle touches the envelope
\(\tau = c + \sigma_n \tan\phi\). The same two quantities are used in
Section~\ref{sec:hb-to-mc-intact} to fit the envelope to strength data.

\paragraph{Limitations.}
The criterion is linear, and so does not capture the curvature of rock
strength across different levels of confinement. Extended into tension it
often overestimates the tensile strength. It takes no explicit account of the
intermediate principal stress \(\sigma_2\). And the parameters it provides are
equivalent values that hold only over the stress range from which they were
derived. Mohr--Coulomb is convenient, but it is not a universal material law,
and it is best treated as a design approximation that depends on the stress
range.

\subsection{The Hoek--Brown criterion}
\label{sec:hoek-brown-intact}

\paragraph{The criterion for intact rock.}
For intact rock the empirical Hoek--Brown
criterion~\citep{hoekbrown1980underground,hoek2019hb} is
\[
\sigma_1^f =
\sigma_3 + \sigma_{ci}
\left(
m_i\frac{\sigma_3}{\sigma_{ci}} + 1
\right)^{1/2} ,
\]
or equivalently
\[
\sigma_1^f =
\sigma_3 +
\left(
m_i\sigma_{ci}\sigma_3 + \sigma_{ci}^{2}
\right)^{1/2} ,
\]
where \(\sigma_{ci}\) is the uniaxial compressive strength of the intact
rock, \(m_i\) a material constant, \(\sigma_3\) the confining stress and
\(\sigma_1^f\) the major principal stress at failure. The criterion captures
the nonlinear increase of rock strength with confinement, which is why it is
widely used to derive strength parameters in rock engineering.

\paragraph{The material constant \(m_i\).}
The constant \(m_i\) sets the curvature of the envelope and depends on rock
type. Table~\ref{tab:mi_values} gives typical values for common rock types;
those not listed there are in Table~\ref{tab:mi} of
Appendix~\ref{app:constants}. Where a project supplies its own \(m_i\), that
value takes precedence over the tabulated one. The four rocks of the data set
used in Section~\ref{sec:design-data-and-stress-calculation} carry
\[
\text{Granite}=32,\quad
\text{Sandstone}=17,\quad
\text{Mudstone}=4,\quad
\text{Limestone}=12 ,
\]
and a value once adopted is changed only with explicit justification.

\begin{table}[htbp]
\caption{Typical \(m_i\) values for intact rock.}
\label{tab:mi_values}
\centering
\small
\begin{tabularx}{0.75\textwidth}{l c}
\toprule
\textbf{Rock type} & \textbf{\(m_i\) range} \\
\midrule
Granite & \(32 \pm 3\) \\
Sandstone & \(17 \pm 4\) \\
Siltstone & \(7 \pm 2\) \\
Mudstone & \(4 \pm 2\) \\
Shale & \(6 \pm 2\) \\
Limestone, crystalline & \(12 \pm 3\) \\
Marble & \(9 \pm 3\) \\
Quartzite & \(20 \pm 3\) \\
\bottomrule
\end{tabularx}
\end{table}

\paragraph{Tensile strength.}
Setting \(\sigma_1^f = 0\) and solving for \(\sigma_3\) gives
\[
\sigma_t =
\frac{\sigma_{ci}}{2}
\left(
m_i - \sqrt{m_i^2+4}
\right) .
\]
This is generally lower, and more realistic, than the tensile strength
predicted by Mohr--Coulomb.

\paragraph{Worked example: marble.}
For marble with \(\sigma_{ci}=125~\text{MPa}\) and \(m_i=9\) under a
confining stress \(\sigma_3=20~\text{MPa}\),
\[
\sigma_1^f =
20 + 125
\left(
9\cdot\frac{20}{125}+1
\right)^{1/2}
=
20 + 125(2.44)^{1/2}
=
215~\text{MPa} ,
\]
and the tensile strength is
\[
\sigma_t =
\frac{125}{2}
\left(
9-\sqrt{81+4}
\right)
=
62.5(9-9.22)
=
-13.7~\text{MPa} .
\]

\subsection{Griffith theory and the tensile cut-off}
\label{sec:griffith-tensile-cutoff}

\paragraph{The Griffith criterion.}
Griffith theory considers failure that originates at pre-existing
cracks~\citep{griffith1921rupture}. For compressive stress states it gives
\[
(\sigma_1-\sigma_3)^2 =
8|\sigma_t|(\sigma_1+\sigma_3) ,
\]
which for \(\sigma_3=0\) reduces to
\[
\sigma_c = 8|\sigma_t| .
\]
The corresponding Mohr envelope is parabolic,
\[
\tau^2 =
4|\sigma_t|(\sigma_n+|\sigma_t|) .
\]
The value of the theory is conceptual. It explains why the tensile strength
of a brittle material such as rock is much lower than its compressive
strength, although the ratio of eight that it predicts is lower than those
ordinarily measured.

\paragraph{The Hoek--Martin tensile cut-off.}
\citet{hoekmartin2014fracture} proposed the empirical relationship
\[
\frac{\sigma_{ci}}{|\sigma_t|}
=
8.62 + 0.7m_i ,
\]
which rearranges to
\[
\sigma_t =
-\frac{\sigma_{ci}}{8.62+0.7m_i} .
\]
For the marble of the preceding example the ratio is 14.9 and the tensile
strength \(-8.4~\text{MPa}\), about three-fifths of the value the Hoek--Brown
expression returns. The tensile strength of rock is very low in comparison
with its compressive strength, and this governs the interpretation of crack
initiation, spalling and brittle failure around underground openings
(Chapter~\ref{ch:high-stress}).

\section{From criterion to design parameters}
\label{sec:criteria-in-application}

\subsection{Stress demand and strength capacity}
\label{sec:demand-and-capacity}

The symbols \(\sigma_1\) and \(\sigma_3\) are used in two related but
different senses, and the difference is a common source of error.

\paragraph{Two meanings of the principal stresses.}
In the in-situ context, \(\sigma_1^{IS}\), \(\sigma_2^{IS}\) and
\(\sigma_3^{IS}\) are the principal stresses already acting in the ground
before excavation, obtained by ranking \(\sigma_v\), \(\sigma_H\) and
\(\sigma_h\):
\[
\sigma_1^{IS} = \max(\sigma_v,\sigma_H,\sigma_h) ,
\qquad
\sigma_3^{IS} = \min(\sigma_v,\sigma_H,\sigma_h) .
\]
They are the stress demand imposed by the ground. In the Hoek--Brown
criterion,
\[
\sigma_1^f =
\sigma_3 + \sigma_{ci}
\left(
m_i\frac{\sigma_3}{\sigma_{ci}} + 1
\right)^{1/2} ,
\]
\(\sigma_3\) is the confining stress applied to the specimen or the rock
mass and \(\sigma_1^f\) the major principal stress at failure. These describe
the strength capacity of the material and are not the in-situ stresses. In
general
\[
\sigma_1^f \neq \sigma_1^{IS} ,
\]
since \(\sigma_1^{IS}\) is the maximum stress now in the ground and
\(\sigma_1^f\) the stress required to cause failure at a given confinement.
Failure is assessed by comparing the demand with the capacity, and
Table~\ref{tab:notation_insitu_criterion} gives the notation that keeps the
two apart.

\begin{table}[htbp]
\caption{Notation distinguishing in-situ stresses from the stresses of a failure criterion.}
\label{tab:notation_insitu_criterion}
\centering
\small
\begin{tabularx}{\textwidth}{l l Y}
\toprule
\textbf{Symbol} & \textbf{Meaning} & \textbf{Use} \\
\midrule
\(\sigma_v\) & Vertical in-situ stress & Calculated from \(\gamma z\). \\
\(\sigma_H\) & Maximum horizontal in-situ stress & Calculated from \(k_H\sigma_v\). \\
\(\sigma_h\) & Minimum horizontal in-situ stress & Calculated from \(k_h\sigma_v\). \\
\(\sigma_1^{IS}\) & Maximum in-situ principal stress & Current stress demand in the ground. \\
\(\sigma_3^{IS}\) & Minimum in-situ principal stress & Current minimum confinement in the ground. \\
\(\sigma_3\) in HB & Confining stress used in the failure criterion & Input variable for the HB strength envelope. \\
\(\sigma_1^f\) in HB & Major principal stress at failure & Strength capacity predicted by HB. \\
\(\sigma_{3,\mathrm{fit,max}}\) & Upper confining stress used for HB--MC fitting & Modelling choice used to define the regression range. \\
\bottomrule
\end{tabularx}
\end{table}

\paragraph{The fitting limit \(\sigma_{3,\mathrm{fit,max}}\).}
The value \(\sigma_{3,\mathrm{fit,max}}\) is the upper limit of the range of
confining stress over which the nonlinear Hoek--Brown envelope is replaced by
an equivalent linear Mohr--Coulomb envelope. For preliminary design the range
is taken as
\[
0 \leq \sigma_3 \leq \sigma_{3,\mathrm{fit,max}} ,
\qquad
\sigma_{3,\mathrm{fit,max}}=\max(\sigma_v,\sigma_{h,k}) ,
\]
which covers the main stress level of the project in the simplified
two-stress description. The minimum in-situ principal stress is still
\[
\sigma_3^{IS}=\min(\sigma_v,\sigma_{h,k}) ,
\]
and the two coincide only where \(\sigma_v=\sigma_{h,k}\). The first is a
stress state in the ground. The second is a modelling choice made for curve
fitting.

Fitting to \(\sigma_3^{IS}\) instead is possible, but it gives a narrower
range. A narrower range usually gives a higher \(\phi\) and a lower \(c\),
and a wider range a lower \(\phi\) and a higher \(c\). A fitted pair is
therefore meaningful only with its fitting range stated and justified.

\paragraph{The sequence of calculation.}
The calculation proceeds in a fixed order. The vertical stress
\(\sigma_v = \gamma z\) is found first, and the horizontal stress
\(\sigma_{h,k} = k\sigma_v\) from the stress ratio supplied for the project,
where \(\sigma_{h,k}\) is the horizontal stress of the simplified
two-dimensional calculation. The two are ranked to give \(\sigma_1^{IS}\) and
\(\sigma_3^{IS}\). The intact strength is then evaluated from the Hoek--Brown
criterion at the confinement of each rock, \(\sigma_3 = \sigma_3^{IS}\), and
compared with \(\sigma_1^{IS}\). For the Mohr--Coulomb conversion,
Hoek--Brown data are generated over the range
\(0 \leq \sigma_3 \leq \sigma_{3,\mathrm{fit,max}}\). Finally, the fitting
range is stated together with the reason it is relevant to the stress
condition of the project.

A design report records the same logic in the same order: the in-situ
stresses calculated from depth and stress ratio define the stress condition
of the project and represent the demand; the Hoek--Brown criterion gives the
capacity of the intact rock; and the equivalent Mohr--Coulomb parameters are
obtained by fitting the Hoek--Brown envelope over the confinement range
relevant to the project.

\subsection{Equivalent Mohr--Coulomb parameters}
\label{sec:hb-to-mc-intact}

Because the Hoek--Brown envelope is nonlinear, equivalent Mohr--Coulomb
parameters \(c\) and \(\phi\) are obtained by fitting a straight line to
Hoek--Brown data over a selected stress range. They are not unique: they
depend on the range, which is why the limit
\(\sigma_{3,\mathrm{fit,max}}\) of the preceding section is fixed before any
fitting is done.

\paragraph{Procedure.}
The intact parameters \(\sigma_{ci}\) and \(m_i\) are selected, the in-situ
stresses \(\sigma_v\) and \(\sigma_{h,k}\) calculated, and the fitting range
defined from them. Several \((\sigma_3,\sigma_1^f)\) pairs are then generated
from the Hoek--Brown equation across the range. The pairs are either
converted to Mohr circle centres and radii or fitted directly as
\(\sigma_1^f\) against \(\sigma_3\), and \(c\) and \(\phi\) follow from the
coefficients of the fitted line. The last step is to check that the line
represents the curved envelope reasonably over the chosen range.

\paragraph{Method 1: Mohr circle regression.}
Each \((\sigma_3,\sigma_1^f)\) pair defines a circle with
\[
C = \frac{\sigma_1^f+\sigma_3}{2},
\qquad
R = \frac{\sigma_1^f-\sigma_3}{2} .
\]
Linear regression on the \((C,R)\) pairs,
\[
R = A + B C ,
\]
gives
\[
\phi = \arcsin(B),
\qquad
c = \frac{A}{\cos\phi} .
\]
The gradient $B$ and the intercept $A$ are those of an ordinary
least-squares fit, available in any
spreadsheet or statistical package and implemented in \companionname\
(Appendix~\ref{app:software}). What matters is not the tool but the range:
the fit must be made over the confining stresses the problem actually spans,
because the equivalent cohesion and friction angle are properties of that
range and not of the rock alone.

\paragraph{Method 2: linear regression in \(\sigma_1^f\)--\(\sigma_3\) space.}
The same data may be fitted directly as
\[
\sigma_1^f = a_{\mathrm{fit}} + b_{\mathrm{fit}}\,\sigma_3 ,
\]
in which the intercept \(a_{\mathrm{fit}}\) is the uniaxial compressive
strength implied by the fitted line and the gradient \(b_{\mathrm{fit}}\)
corresponds to \(k_\phi\). The Mohr--Coulomb parameters are then
\[
\sin\phi = \frac{b_{\mathrm{fit}}-1}{b_{\mathrm{fit}}+1} ,
\qquad
c =
\frac{a_{\mathrm{fit}}(1-\sin\phi)}
{2\cos\phi} .
\]
The two methods give similar results when the same stress range and data
points are used. A large difference usually indicates an arithmetic or unit
error and not a genuine feature of the data.

\subsection{Intact strength of four rocks}
\label{sec:design-data-and-stress-calculation}

\paragraph{Input data.}
Table~\ref{tab:intact_project_data} summarises the intact properties and
project conditions of the four rocks introduced in
Chapter~\ref{ch:classification}.

\begin{table}[htbp]
\caption{Intact properties and project conditions of the four rocks.}
\label{tab:intact_project_data}
\centering
\small
\begin{tabularx}{\textwidth}{l c c c c}
\toprule
\textbf{Parameter} & \textbf{Granite} & \textbf{Sandstone} & \textbf{Mudstone} & \textbf{Limestone} \\
\midrule
\(\sigma_{ci}\) (MPa) & 265 & 110 & 27 & 85 \\
\(E_i\) (GPa) & 80 & 30 & 0.76 & 20 \\
\(I_{s(50)}\) (MPa) & 12.2 & 5.3 & 2.3 & 5.0 \\
\(\sigma_t\) from lab (MPa) & \(-18.4\) & \(-6.7\) & \(-3.1\) & \(-4.5\) \\
RQD (\%) & 70 & 68 & 38 & 52 \\
\(m_i\) & 32 & 17 & 4 & 12 \\
\midrule
Depth \(z\) (m) & 780 & 450 & 100 & 80 \\
Stress ratio \(k\) & 3.2 & 1.8 & 1.0 & 2.0 \\
\(\gamma\) (kN/m\textsuperscript{3}) & 27 & 27 & 27 & 27 \\
\bottomrule
\end{tabularx}
\end{table}

The intact modulus \(E_i\) is not used in this chapter; it is carried forward
to the deformation modulus of the rock mass in
Chapter~\ref{ch:rock-mass-parameters}. Its ratio to the intact strength, the
modulus ratio of Appendix~\ref{app:constants}, is \(80\,000/265 = 302\),
\(30\,000/110 = 273\) and \(20\,000/85 = 235\) for the granite, the sandstone
and the limestone, within the ranges usual for those rocks, but
\(760/27 = 28\) for the mudstone. That is an order of magnitude below the
range of roughly 200 to 400 usual for mudstone, and it implies a
linear-elastic strain at peak of \(27/760\), about 3.6 per cent. The mudstone
value is retained as recorded, but it is out of family with its strength, and
it should be confirmed by testing before a modulus derived from it is relied
on.

\paragraph{In-situ stresses.}
With \(\gamma = 0.027~\text{MPa/m}\),
\[
\sigma_v = \gamma z,
\qquad
\sigma_{h,k} = k\sigma_v .
\]
Only one horizontal stress ratio is available for each rock, so the
two-stress comparison applies, with
\(\sigma_1^{IS}=\max(\sigma_v,\sigma_{h,k})\),
\(\sigma_3^{IS}=\min(\sigma_v,\sigma_{h,k})\) and the fitting limit
\(\sigma_{3,\mathrm{fit,max}}=\max(\sigma_v,\sigma_{h,k})\).
Table~\ref{tab:insitu_stress_fitting} gives the results. The horizontal
stress is the major principal stress in the granite, the sandstone and the
limestone, and the mudstone, with \(k = 1.0\), is the one case in which the
minimum in-situ stress and the fitting limit coincide.

\begin{table}[htbp]
\caption{In-situ principal stresses of the four rocks and the upper fitting limit for the conversion to Mohr--Coulomb parameters. Stresses in MPa.}
\label{tab:insitu_stress_fitting}
\centering
\small
\begin{tabular}{l c c c c c c}
\toprule
\textbf{Rock} & \(\sigma_v\) & \(k\) & \(\sigma_{h,k}\) & \(\sigma_1^{IS}\) & \(\sigma_3^{IS}\) & \(\sigma_{3,\mathrm{fit,max}}\) \\
 & (MPa) & & (MPa) & (MPa) & (MPa) & (MPa) \\
\midrule
Granite & 21.06 & 3.2 & 67.39 & 67.39 & 21.06 & 67.39 \\
Sandstone & 12.15 & 1.8 & 21.87 & 21.87 & 12.15 & 21.87 \\
Mudstone & 2.70 & 1.0 & 2.70 & 2.70 & 2.70 & 2.70 \\
Limestone & 2.16 & 2.0 & 4.32 & 4.32 & 2.16 & 4.32 \\
\bottomrule
\end{tabular}
\end{table}

\paragraph{Strength at in-situ confinement.}
Table~\ref{tab:intact_strength_results} evaluates the intact Hoek--Brown
criterion for each rock at its own \(\sigma_3^{IS}\). In every case the
intact strength exceeds the major in-situ stress several times over, by a
factor of about 8 in the granite and 9 in the sandstone. That margin belongs
to the intact material. The strength that governs an excavation is that of
the jointed rock mass, which Chapter~\ref{ch:rock-mass-parameters} derives
from the same data and which is very much lower.

\begin{table}[htbp]
\caption{Intact strength of the four rocks at in-situ confinement, and
tensile strength from the Hoek--Brown expression, from the Hoek--Martin
cut-off and from laboratory tests. Stresses in MPa.}
\label{tab:intact_strength_results}
\centering
\small
\begin{tabular}{l c c c c c c}
\toprule
\textbf{Rock} & \(\sigma_3^{IS}\) & \(\sigma_1^f\) & \(\sigma_1^{IS}\)
 & \(\sigma_t\), Hoek--Brown & \(\sigma_t\), Hoek--Martin & \(\sigma_t\), laboratory \\
\midrule
Granite & 21.06 & 519.9 & 67.39 & \(-8.3\) & \(-8.5\) & \(-18.4\) \\
Sandstone & 12.15 & 198.8 & 21.87 & \(-6.5\) & \(-5.4\) & \(-6.7\) \\
Mudstone & 2.70 & 34.6 & 2.70 & \(-6.4\) & \(-2.4\) & \(-3.1\) \\
Limestone & 2.16 & 99.3 & 4.32 & \(-7.0\) & \(-5.0\) & \(-4.5\) \\
\bottomrule
\end{tabular}
\end{table}

The three estimates of tensile strength do not agree. For the sandstone all
three lie within about 20 per cent of one another. For the mudstone, with its
low \(m_i\), the Hoek--Brown expression returns twice the measured value and
the cut-off is the closer estimate, whereas for the granite both give less
than half the laboratory magnitude. Where tensile failure matters, the measured
value governs and the two expressions serve as checks on it.

\paragraph{Equivalent parameters for the sandstone.}
The sandstone, with \(\sigma_{ci}=110~\text{MPa}\) and \(m_i=17\), is fitted
over \(0 \leq \sigma_3 \leq 21.87~\text{MPa}\).
Table~\ref{tab:sandstone_hb_data} lists six evenly spaced points generated
from the intact criterion.

\begin{table}[htbp]
\caption{Values of \(\sigma_1^f\) generated for the sandstone from the intact
Hoek--Brown criterion over its fitting range.}
\label{tab:sandstone_hb_data}
\centering
\small
\begin{tabular}{c c c c}
\toprule
\(\sigma_3\) (MPa) & \(\sigma_1^f\) (MPa) & Centre \(C\) (MPa) & Radius \(R\) (MPa) \\
\midrule
0 & 110.00 & 55.00 & 55.00 \\
4.37 & 146.78 & 75.58 & 71.20 \\
8.75 & 177.45 & 93.10 & 84.35 \\
13.12 & 204.53 & 108.83 & 95.71 \\
17.50 & 229.20 & 123.35 & 105.85 \\
21.87 & 252.08 & 136.98 & 115.11 \\
\bottomrule
\end{tabular}
\end{table}

Method 1 gives \(R = 15.50 + 0.7325\,C\), from which
\(\phi = 47.1^\circ\) and \(c = 22.8~\text{MPa}\). Method 2 gives
\(\sigma_1^f = 116.3 + 6.432\,\sigma_3\), from which \(\phi = 47.0^\circ\)
and \(c = 22.9~\text{MPa}\). The two agree, as they should on the same
points. The intercept of the second line, 116.3~MPa, exceeds
\(\sigma_{ci}\) by about 6 per cent, and the tensile strength the fitted pair
implies is \(-18.1~\text{MPa}\) against a measured \(-6.7~\text{MPa}\): a
straight line cannot follow the curved envelope at both ends, and it departs
furthest in tension.

Fitting over the narrower range \(0 \leq \sigma_3 \leq 12.15~\text{MPa}\),
the minimum in-situ stress, gives \(\phi = 49.3^\circ\) and
\(c = 20.9~\text{MPa}\) by the second method, a higher friction angle and a
lower cohesion as expected. Neither pair is the more correct. Each describes
the envelope over its own range, and a pair quoted without that range cannot
be used safely. Chapter~\ref{ch:rock-mass-parameters} repeats the conversion
for the same sandstone as an undisturbed rock mass of GSI~63, where the
friction angle falls by
about five degrees and the cohesion to less than a quarter of its intact
value.

\chapterreferences
